\section{Resultados, insights y conclusiones}
\label{sec:resultados}

\subsection{Patrones observados}
\begin{enumerate}[leftmargin=1.5em]
  \item \textbf{La Fase 1 resuelve (o casi) muchas instancias.} En la mayoría de OR-Library
        $LB_1=UB_1=$ óptimo (Tabla~\ref{tab:orlib}): la relajación lineal del maestro de Benders es
        muy ajustada y el redondeo encuentra la solución óptima. La Fase 2 solo \emph{certifica}.
        Coincide con lo reportado por el paper (su Fase 1 ya deja muchas instancias resueltas).
  \item \textbf{El método es correcto en la escala local probada.} La evidencia consolidada
        verifica OR-Library \texttt{pmed1}--\texttt{pmed40} contra óptimos oficiales
        (Tabla~\ref{tab:orlib40}) y 9/9 óptimos en \texttt{rl1304} ($N=1304$, hasta $p=500$) con
        $\Delta=0$ respecto del paper (Tabla~\ref{tab:paper}). El separador, validado por tres
        vías (mano, test, oráculo), es la base de esta confianza. Esta es evidencia de replicación
        computacional parcial, no de la campaña completa del paper.
  \item \textbf{El warm-start reduce el árbol, no necesariamente el tiempo.} Heredar los cortes de
        Fase 1 baja los nodos de B\&B de cientos a $1$--$7$ (Tabla~\ref{tab:warm}), pero en
        instancias sub-segundo la precarga de cientos de cortes puede costar más que el ahorro: el
        beneficio en wall-time se espera en instancias grandes, no medidas aquí.
  \item \textbf{Sin PopStar, el $UB_1$ es peor pero el $LB_1$ coincide.} PopStar es la
        heurística primal externa usada por el paper para construir buenas soluciones iniciales y,
        por tanto, mejores cotas superiores. En \texttt{rl1304} $p\in\{10,100,300\}$ nuestro
        $UB_1$ supera al del paper (p.\,ej. $2243679$ vs $2134295$), mientras el $LB_1$ coincide a
        $\pm1$. La cota inferior depende del Benders implementado; la cota superior depende de la
        heurística primal. Para un benchmark computacional completo contra el paper, PopStar (o
        una heurística equivalente y documentada) debería implementarse o integrarse.
\end{enumerate}

\subsection{Ventajas del enfoque}
\begin{itemize}[noitemsep]
  \item \textbf{Separación en forma cerrada $O(NM)$:} no se resuelve ningún LP por corte; el costo
        por iteración es lineal.
  \item \textbf{Maestro pequeño:} solo $y$ y $\theta$; las variables $z$ (potencialmente
        $\sim N\times M$) nunca se materializan.
  \item \textbf{Solo cortes de optimalidad:} al ser $SP_i$ siempre factible y acotado, no hay
        lógica de factibilidad — implementación más simple y robusta.
  \item \textbf{Cotas tempranas útiles:} $LB_1/UB_1$ de Fase 1 ya acotan fuertemente el óptimo,
        valioso para decidir con tiempo limitado en un contexto real.
\end{itemize}

\subsection{Limitaciones}
\begin{itemize}[noitemsep]
  \item \textbf{Escala probada:} benchmarks con óptimo externo hasta \texttt{rl1304} ($N=1304$) y
        OR-Library hasta \texttt{pmed40}; estrés sintético local hasta $N=5000$ sin timeouts. Los
        benchmarks enormes del paper (TSP $>10^4$, BIRCH, ODM, RW grande) \emph{no} se ejecutaron;
        las afirmaciones sobre ellos son trabajo futuro, no resultados locales.
  \item \textbf{Sin comparación directa con Zebra:} no se reejecutó. Sí es válido mostrar que el
        método local reproduce óptimos del paper y contrastar, como literatura, contra los tiempos
        Zebra reportados por Duran-Mateluna, Ales \& Elloumi (2023). Lo que no es válido es
        presentar una superioridad local frente a Zebra sin ejecutar Zebra bajo un protocolo propio.
  \item \textbf{Mejoras no implementadas:} PopStar, reduced-cost fixing y constraint reduction
        quedan documentadas pero fuera del alcance.
  \item \textbf{Construcción de $S$:} a gran escala dominaría el tiempo total ($O(NM\log M)$); no
        se midió de forma aislada por no haber corrido instancias grandes.
\end{itemize}

\subsection{Aprendizajes para el problema real y para el uso de optimización}
\begin{itemize}[noitemsep]
  \item \textbf{La formulación importa más que el solver.} F1--F4 tienen el mismo óptimo, la misma
        relajación LP y por tanto la misma \emph{brecha de integridad}
        $v(\text{MILP})-v(\text{LP})$; sin embargo su estructura (densidad, número de variables)
        decide qué es resoluble. La ganancia vino de \emph{explotar estructura}, no de fuerza
        bruta del solver. Que muchas instancias cierren ya en Fase 1 ($LB_1=UB_1=$ óptimo) es
        reflejo de una formulación \emph{fuerte}: brecha de integridad casi nula.
  \item \textbf{Las buenas cotas tempranas son accionables.} Para un decisor de localización, un
        $LB_1/UB_1$ ajustado en segundos permite descartar configuraciones sin esperar la prueba
        de optimalidad completa.
  \item \textbf{Verificar el componente crítico paga.} Validar el separador a mano, con test y con
        oráculo dio una base de confianza que permitió escalar sin temor a errores silenciosos.
\end{itemize}

\subsection{Conclusión}
Se replicó e implementó en C, de forma modular y reproducible, la descomposición de Benders
eficiente para el pMP. La contribución local es una \textbf{implementación y verificación completa} del
mecanismo Benders central del paper —F3, Benders, separación cerrada, dos fases y warm-start— más una
\textbf{replicación computacional parcial documentada}. El método alcanza el óptimo en todas las
instancias probadas (OR-Library \texttt{pmed1}--\texttt{pmed40} con óptimos oficiales,
\texttt{rl1304} 9/9 vs paper; catálogo local \texttt{paperbench} 26/26 OK; estrés sintético
$N=5000$ sin timeouts), con un separador verificado por tres vías independientes. No se
reejecutó Zebra ni la campaña completa de gran escala. El trabajo cumple los objetivos del curso:
formulación rigurosa, propuesta lagrangiana, diseño y justificación de Benders, e implementación
con análisis honesto de alcances y límites.
